In their paper, \citet{everitt} show that, for modification-independent belief and utility function, if we start with a perfect expected utility maximizer and at any time replace the current policy by the initial policy, the expected discounted utility stays the same. Therefore, later policies cannot be worse than the initial policy and no deterioration happens. We show that in the case of $\epsilon$-optimizers, such a replacement can never decrease the expected discounted utility by more than $\epsilon$ (Inequality \eqref{eq:1lower_bound_policy_mod}) but can increase it more, meaning that agent's behaviour can deteriorate with time (Inequality \eqref{eq:1upper_bound_policy_mod}). Specifically, it can deteriorate at an exponential rate, until its actions become arbitrarily bad -- that is, until the expected future utility lost is the maximum achievable utility (which is $\frac{1}{1-\gamma}$).
\begin{theorem} \label{thm:policy_modification}
Let $\rho$ and $u$ be modification-independent. Consider a self-modifying agent which is an $\epsilon$-optimizer for the empty history. Then, for every $t \geq 1$, \begin{align}
E_{\mae_{<t}} [Q\left(\mae_{<t} \pi_{t}\left(\mae_{<t}\right)\right)] \geq E_{\mae_{<t}}[Q\left(\mae_{<t} \pi_{1}\left(\mae_{<t}\right)\right)] - \min(\frac{\epsilon}{\gamma^{t-1}}, \frac{1}{1-\gamma}) \label{eq:1upper_bound_policy_mod} %AAAI
%E_{\mae_{<t}} [Q\left(\mae_{<t} \pi_{t}\left(\mae_{<t}\right)\right)] \geq &E_{\mae_{<t}}[Q\left(\mae_{<t} \pi_{1}\left(\mae_{<t}\right)\right)]\\& - \min(\frac{\epsilon}{\gamma^{t-1}}, \frac{1}{1-\gamma}) \label{eq:1upper_bound_policy_mod}
\end{align}
where the expectation is with respect to $\mae_{<t}$ such that the perceptions are distributed according to the belief and the actions are given by $a_{i}=\pi_{i}\left(\mae_{<i}\right)$.

Moreover, for all histories $\mae_{<t}$ given by $a_{i}=\pi_{i}\left(\mae_{<i}\right)$ for which the agent is an $\epsilon$-optimizer it holds that
\begin{gather}
Q\left(\mae_{<t} \pi_{1}\left(\mae_{<t}\right)\right) + \epsilon \geq Q\left(\mae_{<t} \pi_{t}\left(\mae_{<t}\right)\right) \label{eq:1lower_bound_policy_mod}\\
\end{gather}
Equality in Inequality \eqref{eq:1upper_bound_policy_mod} can be achieved up to a factor of at most $\gamma$.
\end{theorem}

The expectation in inequality \eqref{eq:1upper_bound_policy_mod} is necessary as can be demonstrated by the following example. Consider an environment in which the first perception is $\alpha$ with probability $\epsilon (1-\gamma)$ and $\beta$ otherwise. Regardless of the first perception, the utility in the future is always $1$ if the action following this perception is $a$ and $0$ if $b$. An $\epsilon$-optimizing agent which performs action $a$ may choose to self-modify to an agent which performs action $a$ after perception $\alpha$ and $b$ after perception $\beta$, thus losing $\frac{\gamma}{1-\gamma}$ utility in the case of perception $\beta$, regardless of how small $\epsilon$ is.

\smallskip \noindent
Setting $\epsilon=0$ allows us to easily recover Theorem 16 from \cite{everitt}, showing that self-modifications do not impact expected discounted utility gained by perfectly rational agents. This proof is also considerably simpler than the one in the original paper. For the proof, see \Cref{cor:recover_everitt}. %AAAI

If we only care about future discounted utility, this deterioration in performance doesn't need to concern us because it only happens at future times when utility is heavily discounted. From the definition of an $\epsilon$-optimizer, the maximum utility lost is indeed only $\epsilon$. On the other hand, if we care about long-term performance of the agent and have only introduced the discount factor for instrumental reasons (as would likely be the case), the possibility of self-modification becomes a serious problem. The discount factor might be introduced because optimizing the long-term future might be computationally intractable.
